\documentclass[12pt,a4paper]{article}

\usepackage{amsmath}
\usepackage{amssymb}
\usepackage[margin=1in]{geometry}

\title{Mathematical Equations in LaTeX}
\author{Abhishek Narayan Hathwala \\ MNIT Jaipur \\ Student ID: 2025UCP1029}
\date{\today}

\begin{document}

\maketitle

\section{Introduction}

LaTeX provides powerful mathematical typesetting features.
It is widely used to write equations, formulas, and
scientific expressions.

\section{Basic Mathematical Expressions}

Inline equation: $a^2 + b^2 = c^2$.

Display equation:

\[
x^2 + y^2 = z^2
\]

\section{Fractions}

A simple fraction:

\[
\frac{a+b}{c+d}
\]

A complex fraction:

\[
\frac{\frac{a}{b}+\frac{c}{d}}{x+y}
\]

\section{Square Roots}

Square root:

\[
\sqrt{x^2+y^2}
\]

Cube root:

\[
\sqrt[3]{x+y}
\]

\section{Superscripts and Subscripts}

\[
x_1, x_2, x_3, \ldots, x_n
\]

\[
a^2+b^2+c^2
\]

\[
x_i^{n+1}
\]

\section{Greek Letters}

Some commonly used Greek letters are:

\[
\alpha, \beta, \gamma, \delta, \theta, \lambda
\]

\[
\pi, \sigma, \omega, \mu, \phi
\]

\section{Summation and Product}

Summation:

\[
\sum_{i=1}^{n} i = \frac{n(n+1)}{2}
\]

Product:

\[
\prod_{i=1}^{n} i = n!
\]

\section{Limits}

The standard limit expression:

\[
\lim_{x \to 0} \frac{\sin x}{x} = 1
\]

\section{Integration}

Definite integration:

\[
\int_0^1 x^2\,dx = \frac{1}{3}
\]

Indefinite integration:

\[
\int x^2\,dx = \frac{x^3}{3}+C
\]

\section{Aligned Equations}

The align environment is used to align multiple equations.

\begin{align}
a+b &= c \\
x+y &= z \\
m+n &= p
\end{align}

\section{Conclusion}

This practical demonstrates how LaTeX can be used to
create mathematical expressions, fractions, roots,
Greek symbols, summations, limits, and integrals.

\end{document}